\documentclass[../../main.tex]{subfiles}

\graphicspath{{../../image/}{../../figure/}}

\begin{document}

\section{偶然熵与多凸性}\label{section:ch05-偶然熵与多凸性}

我们在此考虑带对合与非凸偶然熵的守恒律方程组，如前一节那样，但还装备有补充偶然熵，使 Cauchy 问题解的存在性与稳定性像\cref{theorem:ch05-凸熵局部存在}所推断的那样强。

弹性动力学方程组\eqref{align:ch03-等熵热弹性系统}与电动力学方程组\eqref{align:ch03-Maxwell}将作为模型。

如前一节所指出，\eqref{align:ch03-等熵热弹性系统}装备有对合\eqref{align:ch03-协调方程}、

熵 $\eta=\varepsilon(F)+\frac{1}{2}|v|^{2}$ 及十个偶然熵，即 $\det F$ 与伴随矩阵 $F^{*}$ 的九个元素。类似地，\eqref{align:ch03-Maxwell}装备有对合\eqref{align:ch03-高斯}、

电磁场能量 $\eta(B,D)$ 及三个偶然熵，即向量 $B\wedge D$ 的分量。因此，贯穿本节我们将在如下假设下工作。

\begin{assumption}[多凸性基本假设]\label{assumption:ch05-polyconvex-basic}
在守恒律方程组\eqref{align:ch05-cauchy-system}中，通量 $G$ 满足对称条件\eqref{align:ch05-symmetry-condition}，它诱导对合\eqref{align:ch05-involution}。方程组装备有\emph{主偶然熵--熵通量对} $(\eta,Q)$ 及一族 $N$ 个补充偶然熵--熵通量对 $(W^{I},X^{I})$，$I=1,\ldots,N$。
\end{assumption}

作为偶然熵对，$(\eta,Q)$ 须对某 $k$-向量值函数 $\Xi(U)$ 满足\eqref{align:ch05-contingent-def}。类似地，对 $I=1,\ldots,N$，
\begin{align}
\mathrm{D}X_{\alpha}^{I}(U)=\mathrm{D}W^{I}(U)\mathrm{D}G_{\alpha}(U)+\Omega^{I}(U)^{\top}M_{\alpha},\qquad\alpha=1,\ldots,m,\label{align:ch05-supplementary}
\end{align}
对某 $k$-向量值函数 $\Omega^{I}(U)$。将 $W^{I}$ 组装为 $N$-向量 $W$、$X^{I}$ 组装为 $N\times m$ 矩阵 $X$、$\Omega^{I}$ 组装为 $N\times k$ 矩阵 $\Omega$ 是方便的，此时\eqref{align:ch05-supplementary}可写成矩阵方程
\begin{align}
\mathrm{D}X_{\alpha}(U)=\mathrm{D}W(U)\mathrm{D}G_{\alpha}(U)+\Omega(U)^{\top}M_{\alpha},\qquad\alpha=1,\ldots,m.\label{align:ch05-supplementary-matrix}
\end{align}

状态向量 $U$ 的每个分量 $U^{i}$ 可视为一个熵，相应的熵为矩阵 $G$ 的第 $i$ 行 $G^{i}$。为方便起见，我们将这些熵对嵌入\cref{assumption:ch05-polyconvex-basic}记录的补充对列表中。于是我们假设 $N\geqslant n$，且对 $I=1,\ldots,n$，$W^{I}(U)=U^{I}$、$X_{\alpha}^{I}(U)=G_{\alpha}^{I}(U)$、$\Omega^{I}(U)=0$。

方程\eqref{align:ch05-entropy-conservation}连同
\begin{align}
\partial_{t}W(U(x,t))+\sum_{\alpha=1}^{m}\partial_{\alpha}X_{\alpha}(U(x,t))=0\label{align:ch05-W-conservation}
\end{align}
须对满足对合\eqref{align:ch05-involution}的\eqref{align:ch05-cauchy-system}的任一经典解 $U$ 成立。另一方面，可容许弱解 $U$ 须对主偶然熵--熵通量对 $(\eta,Q)$ 满足不等式\eqref{align:ch04-熵积分0}。
典型地，在连续介质物理中遇到的带对合的方程组（如\eqref{align:ch03-等熵热弹性系统}与\eqref{align:ch03-Maxwell}）中，

主偶然熵是真正的熵。

本节的目的是证明：在上述设定下，主熵上的凸性要求可放宽为如下更弱条件。

\begin{definition}[多凸]\label{definition:ch05-polyconvex}
主偶然熵 $\eta$ 称为相对于偶然熵 $W$ \emph{多凸}（polyconvex），若它可表示为
\begin{align}
\eta(U)=\theta(W(U)),\qquad U\in\overline{\mathscr{B}_{\rho}},\label{align:ch05-polyconvex-def}
\end{align}
其中 $\theta$ 是定义在 $W(\mathscr{B}_{\rho})$ 在 $\mathbb{R}^{N}$ 中某开邻域 $\mathscr{F}$ 上的光滑函数，其 Hessian 矩阵在每个 $W\in\mathscr{F}$ 上正定。
\end{definition}

在弹性动力学的例子中，$W=(F,v,F^{*},\det F)$ 排列为 22-向量，当内能函数 $\varepsilon(F)$ 容许表示
\begin{align}
\varepsilon(F)=\theta(F,F^{*},\det F),\label{align:ch05-elastic-polyconvex}
\end{align}
时主熵 $\eta=\varepsilon(F)+\frac{1}{2}|v|^{2}$ 将是多凸的，其中 $\theta(F,H,\delta)$ 是在流形 $\{(F,H,\delta):\det F>0,\ H=F^{*},\ \delta=\det F\}\subset\mathbb{R}^{19}$ 的某开邻域上 Hessian 正定的光滑函数。这是一个物理上合理的假设，已在文献中——特别在弹性静力学的背景下——被彻底讨论过。我们已在\cref{section:ch04-熵可容性条件}中遇到它，那里注意到它蕴涵 $\varepsilon(F)$ 秩一凸，从而 $\eta$ 在对合锥上凸。
然而正如我们在此将看到的，多凸性对解稳定性的蕴含比上一节讨论的对合锥上单纯凸性的后果强得多。

电动力学方程组\eqref{align:ch03-Maxwell}的情形类似，此时 $W=(B,D,B\wedge D)$ 排列为 9-向量。
多凸性是电磁场能量 $\eta$（充当主熵）的自然条件。事实上，在 Born--Infeld 情形\eqref{align:ch03-BornInfeld}中，$\eta$ 是多凸的。

我们为函数 $\theta(W)$ 引入如下记号：$\theta_{W^{I}}$ 表示偏导数 $\partial\theta/\partial W^{I}$；$\theta_{W}$ 表示微分 $[\theta_{W^{1}}\cdots\theta_{W^{N}}]$，视为 $1\times n$ 矩阵；$\theta_{WW}$ 表示 $N\times N$ Hessian 矩阵。

对 $U\in\overline{\mathscr{B}_{\rho}}$，我们定义 $n\times n$ 矩阵
\begin{align}
A(U)=\mathrm{D}^{2}\eta(U)-\sum_{I=1}^{N}\theta_{W^{I}}(W(U))\mathrm{D}^{2}W^{I},\label{align:ch05-A-polyconvex}
\end{align}
\begin{align}
J_{\alpha}(U)=A(U)\mathrm{D}G_{\alpha}(U)+\Gamma(U)^{\top}M_{\alpha},\qquad\alpha=1,\ldots,m,\label{align:ch05-Jalpha-polyconvex}
\end{align}
其中
\begin{align}
\Gamma(U)=\mathrm{D}\Xi(U)-\sum_{I=1}^{N}\theta_{W^{I}}(W(U))\mathrm{D}\Omega^{I}(U).\label{align:ch05-Gamma-def}
\end{align}

显然 $A(U)$ 对称。如\cref{section:ch05-对合与偶然熵}已指出，\eqref{align:ch05-contingent-def}蕴涵矩阵
\begin{align}
\mathrm{D}^{2}\eta(U)\mathrm{D}G_{\alpha}(U)+\mathrm{D}\Xi(U)^{\top}M_{\alpha},\qquad\alpha=1,\ldots,m,\label{align:ch05-sym1}
\end{align}
对称。类似地，\eqref{align:ch05-supplementary}蕴涵矩阵
\begin{align}
\mathrm{D}^{2}W^{I}(U)\mathrm{D}G_{\alpha}(U)+\mathrm{D}\Omega^{I}(U)^{\top}M_{\alpha},\qquad\alpha=1,\ldots,m,\label{align:ch05-sym2}
\end{align}
亦对称。于是矩阵 $J_{\alpha}(U)$（$\alpha=1,\ldots,m$）对称。

由\eqref{align:ch05-elastic-polyconvex}与\eqref{align:ch05-A-polyconvex}，
\begin{align}
A(U)=\mathrm{D}W(U)^{\top}\theta_{WW}(W(U))\mathrm{D}W(U),\label{align:ch05-A-factor}
\end{align}
于是若 $\eta$ 多凸，则 $A(U)$ 正定。

下述命题确立了带多凸熵的方程组 Cauchy 问题经典解的局部存在性。

\begin{theorem}\label{theorem:ch05-polyconvex-existence}
设守恒律双曲方程组\eqref{align:ch05-cauchy-system}满足\cref{assumption:ch05-polyconvex-basic}，主偶然熵 $\eta$ 多凸\eqref{align:ch05-polyconvex-def}。设初值 $U_{0}\in H_{\ell}$，某 $\ell>\frac{m}{2}+1$，取值于半径 $\rho_{0}<\rho$ 的球 $\overline{\mathscr{B}_{\rho_{0}}}$，并满足对合\eqref{align:ch05-involution}。则存在 $T_{\infty}\leqslant\infty$ 及 $\mathbb{R}^{m}\times[0,T_{\infty})$ 上唯一连续可微、取值于 $\overline{\mathscr{B}_{\rho}}$ 的函数 $U$，它是 Cauchy 问题\eqref{align:ch05-cauchy-system}、\eqref{align:ch05-cauchy-data}在时间区间 $[0,T_{\infty})$ 上的经典解。此外，
\begin{align}
U(\cdot,t)\in\bigcap_{k=0}^{\ell}C^{k}([0,T_{\infty});H_{\ell-k}).\label{align:ch05-polyconvex-regularity}
\end{align}
区间 $[0,T_{\infty})$ 是极大的，即若 $T_{\infty}<\infty$，则
\begin{align}
\int_{0}^{T_{\infty}}\|\|\nabla U(\cdot,t)\|\|_{L^{\infty}}\,\mathrm{d}t=\infty,\label{align:ch05-polyconvex-blowup1}
\end{align}
且/或
\begin{align}
\limsup_{t\to T_{\infty}}\|U(\cdot,t)\|_{L^{\infty}}=\rho.\label{align:ch05-polyconvex-blowup2}
\end{align}
\end{theorem}

\begin{proof}
遵循粘性消去法，我们将\eqref{align:ch05-cauchy-system}、\eqref{align:ch05-cauchy-data}的解构造为\eqref{align:ch05-parabolic}、\eqref{align:ch05-cauchy-data}的解当 $\varepsilon\to0$ 时的极限，回溯\cref{theorem:ch05-凸熵局部存在}证明中经由引理 5.1.2、5.1.3、5.1.4 的步骤。
事实上，此处记号的设计使得：一旦将\eqref{align:ch05-A-polyconvex}与\eqref{align:ch05-Jalpha-polyconvex}取代\eqref{align:ch05-A-def}与\eqref{align:ch05-Jalpha-def}作为 $A(U)$ 与 $J_{\alpha}(U)$ 的定义，就可将\cref{section:ch05-凸熵与经典解的存在性}的文本与方程几乎逐字搬到当前设定。
直截了当的验证留待读者。
\end{proof}

现转向经典解在将称为\emph{mild}（温和）的弱解类中的唯一性与稳定性问题。

\begin{definition}[mild 解]\label{definition:ch05-mild}
$\mathbb{R}^{m}\times[0,T)$ 上取值于 $\overline{\mathscr{B}_{\rho}}$ 的可测函数 $U$ 称为\eqref{align:ch05-cauchy-system}、\eqref{align:ch05-cauchy-data}的\emph{mild 解}，若对所有在 $\mathbb{R}^{m}\times[0,T)$ 中具紧支集的 Lipschitz $N$-向量值试验函数 $V$，成立
\begin{align}
\int_{0}^{T}\int_{\mathbb{R}^{m}}[\partial_{t}V^{\top}W(U)+\sum_{\alpha=1}^{m}\partial_{\alpha}V^{\top}X_{\alpha}(U)]\,\mathrm{d}x\,\mathrm{d}t+\int_{\mathbb{R}^{m}}V^{\top}(x,0)W(U_{0}(x))\,\mathrm{d}x=0.\label{align:ch05-mild-def}
\end{align}
\end{definition}

注意当 $U$ 在分布意义下满足\eqref{align:ch05-W-conservation}，连同初始条件 $W(U(\cdot,t))\to W(U_{0}(\cdot))$ 于 $L^{\infty}$ 弱$^{*}$、当 $t\to0$ 时，\eqref{align:ch05-mild-def}成立。特别地，\eqref{align:ch05-cauchy-system}、\eqref{align:ch05-cauchy-data}的任一 mild 解都是弱解，因为\eqref{align:ch05-cauchy-system}嵌入\eqref{align:ch05-W-conservation}中。显然，\eqref{align:ch05-cauchy-system}、\eqref{align:ch05-cauchy-data}的任一经典解都是 mild 解，因为此时\eqref{align:ch05-W-conservation}与初始条件自动满足。然而令人惊讶的是，在应用中人们常遇到甚至间断的 mild 解。例如，等熵弹性动力学方程组\eqref{align:ch03-等熵热弹性系统}的任一弱解 $(F,v)$ 是 mild 的。
事实上，如\cref{section:ch02-连续介质热力学平衡律}所见，
\eqref{align:ch03-余因子守恒}与\eqref{align:ch03-行列式守恒}对任一满足\eqref{align:ch03-等熵热弹性系统}$_{1}$及对合\eqref{align:ch03-协调方程}的 $L^{\infty}$ 场成立。



此外，如\cref{section:ch04-熵可容性条件}所注，零 Lagrangian\eqref{align:ch02-2-2-9}是 $L^{\infty}$ 弱$^{*}$ 连续函数，

从而 $F(\cdot,t)\to F_{0}(\cdot)$ 于 $L^{\infty}$ 弱$^{*}$、当 $t\to0$，蕴涵 $F^{*}(\cdot,t)\to F_{0}^{*}(\cdot)$ 与 $\det F(\cdot,t)\to\det F_{0}(\cdot)$ 于 $L^{\infty}$ 弱$^{*}$、当 $t\to0$。类似地，带 Born--Infeld 本构关系\eqref{align:ch03-BornInfeld}与对合\eqref{align:ch03-高斯}的电动力学方程组\eqref{align:ch03-Maxwell}的 $BV$ 弱解 $(B,D)$ 必为 mild 解，


因为所有激波满足\eqref{align:ch03-电磁动量跳跃}。
于是对这些解\eqref{align:ch03-动量电磁}成立。
此外，在 $BV$ 设定下有充分正则性使 $B(\cdot,t)\to B_{0}(\cdot)$ 与 $D(\cdot,t)\to D_{0}(\cdot)$、当 $t\to0$，蕴涵 $B(\cdot,t)\wedge D(\cdot,t)\to B_{0}(\cdot)\wedge D_{0}(\cdot)$、当 $t\to0$。

mild 解 $U$ 称为\emph{可容许}的，若它作为弱解是可容许的，即对主偶然熵--熵通量对满足\eqref{align:ch04-熵积分0}。
特别地，任一 Born--Infeld 本构关系下\eqref{align:ch03-Maxwell}的 $BV$ 解是可容许的，
因为激波不产生能量。当然，弹性动力学方程组\eqref{align:ch03-等熵热弹性系统}则不然。

下述命题是\cref{theorem:ch05-相对熵稳定性}的类比：

\begin{theorem}\label{theorem:ch05-polyconvex-stability}
设守恒律双曲方程组\eqref{align:ch05-cauchy-system}满足\cref{assumption:ch05-polyconvex-basic}，主偶然熵 $\eta$ 多凸\eqref{align:ch05-polyconvex-def}。设 $\bar{U}$ 是\eqref{align:ch05-cauchy-system}在 $[0,T)$ 上的经典解，取值于 $\overline{\mathscr{B}_{\rho}}$，初值 $\bar{U}_{0}$ 满足对合\eqref{align:ch05-involution}。设 $U$ 是\eqref{align:ch05-cauchy-system}在 $[0,T)$ 上的任一可容许 mild 解，取值于 $\overline{\mathscr{B}_{\rho}}$，初值 $U_{0}$ 满足对合\eqref{align:ch05-involution}。则对任意 $r>0$ 与 $t\in[0,T)$，成立
\begin{align}
\int_{|x|<r}|U(x,t)-\bar{U}(x,t)|^{2}\,\mathrm{d}x\leqslant ae^{bt}\int_{|x|<r+st}|U_{0}(x)-\bar{U}_{0}(x)|^{2}\,\mathrm{d}x\label{align:ch05-polyconvex-stability-est}
\end{align}
其中正常数 $s,a$ 仅依赖于 $G,\eta,Q$ 及其在 $\overline{\mathscr{B}_{\rho}}$ 上导数的界，而 $b$ 还依赖于 $\bar{U}$ 的 Lipschitz 常数。特别地，$\bar{U}$ 是以 $\bar{U}_{0}$ 为初值的\eqref{align:ch05-cauchy-system}的唯一可容许 mild 解。
\end{theorem}

\begin{proof}
我们回溯\cref{theorem:ch05-相对熵稳定性}证明的步骤，作必要修改。
在 $\overline{\mathscr{B}_{\rho}}\times\overline{\mathscr{B}_{\rho}}$ 上定义
\begin{align}
h(U,\bar{U})=\eta(U)-\eta(\bar{U})-\theta_{W}(W(\bar{U}))[W(U)-W(\bar{U})],\label{align:ch05-h-polyconvex}
\end{align}
\begin{align}
Y_{\alpha}(U,\bar{U})&=Q_{\alpha}(U)-Q_{\alpha}(\bar{U})-\theta_{W}(W(\bar{U}))[X_{\alpha}(U)-X_{\alpha}(\bar{U})]\nonumber\\
&\quad+[\theta_{W}(W(\bar{U}))\Omega(\bar{U})^{\top}-\Xi(\bar{U})^{\top}]M_{\alpha}[U-\bar{U}],\label{align:ch05-Yalpha-polyconvex}
\end{align}
\begin{align}
Z_{\alpha}(U,\bar{U})&=-\mathrm{D}G_{\alpha}(\bar{U})^{\top}\mathrm{D}W(\bar{U})^{\top}\theta_{WW}(W(\bar{U}))[W(U)-W(\bar{U})]\nonumber\\
&\quad+\mathrm{D}W(\bar{U})^{\top}\theta_{WW}(W(\bar{U}))[X_{\alpha}(U)-X_{\alpha}(\bar{U})]\nonumber\\
&\quad-\mathrm{D}W(\bar{U})^{\top}\theta_{WW}(W(\bar{U}))\Omega(\bar{U})^{\top}M_{\alpha}[U-\bar{U}]\nonumber\\
&\quad+\Gamma(\bar{U})^{\top}M_{\alpha}[U-\bar{U}],\label{align:ch05-Zalpha-polyconvex}
\end{align}
其中 $\Gamma$ 由\eqref{align:ch05-Gamma-def}给出。

回忆\cref{definition:ch05-polyconvex}，我们看到 $h(U,\bar{U})$ 关于 $U-\bar{U}$ 为二阶量且正定。利用\eqref{align:ch05-contingent-def}、\eqref{align:ch05-supplementary-matrix}与\eqref{align:ch05-polyconvex-def}，我们推得
\begin{align}
\mathrm{D}Y_{\alpha}(U,\bar{U})&=[\theta_{W}(W(U))-\theta_{W}(W(\bar{U}))]\mathrm{D}W(U)\mathrm{D}G_{\alpha}(U)\nonumber\\
&\quad+[\Xi(U)-\Xi(\bar{U})]^{\top}M_{\alpha}-\theta_{W}(W(\bar{U}))[\Omega(U)-\Omega(\bar{U})]^{\top}M_{\alpha},\label{align:ch05-DY}
\end{align}
它在 $U=\bar{U}$ 处为零，从而 $Y(U,\bar{U})$ 关于 $U-\bar{U}$ 亦为二阶量。特别地，当 $s$ 大时\eqref{align:ch05-Y-bound}成立。

转向 $Z(U,\bar{U})$，由\eqref{align:ch05-supplementary-matrix}，
\begin{align}
\mathrm{D}Z_{\alpha}(U,\bar{U})&=-\mathrm{D}G_{\alpha}(\bar{U})^{\top}\mathrm{D}W(\bar{U})^{\top}\theta_{WW}(W(\bar{U}))\mathrm{D}W(U)\nonumber\\
&\quad+\mathrm{D}W(\bar{U})^{\top}\theta_{WW}(W(\bar{U}))\mathrm{D}W(U)\mathrm{D}G_{\alpha}(U)\nonumber\\
&\quad+\mathrm{D}W(\bar{U})^{\top}\theta_{WW}(W(\bar{U}))[\Omega(U)-\Omega(\bar{U})]^{\top}M_{\alpha}\nonumber\\
&\quad+\Gamma(\bar{U})^{\top}M_{\alpha}.\label{align:ch05-DZ}
\end{align}

回忆\eqref{align:ch05-A-polyconvex}、\eqref{align:ch05-Jalpha-polyconvex}且 $J_{\alpha}$ 对称，我们结论
\begin{align}
\mathrm{D}Z_{\alpha}(\bar{U},\bar{U})=-\mathrm{D}G_{\alpha}(\bar{U})^{\top}A(\bar{U})+A(\bar{U})\mathrm{D}G_{\alpha}(\bar{U})+\Gamma(\bar{U})^{\top}M_{\alpha}=M_{\alpha}^{\top}\Gamma(\bar{U}).\label{align:ch05-DZ-at-bar}
\end{align}

如\cref{theorem:ch05-相对熵稳定性}证明中那样，固定非负 Lipschitz 连续、在 $\mathbb{R}^{m}\times[0,T)$ 中具紧支集的试验函数 $\psi$，
沿两个解 $U(x,t)$ 与 $\bar{U}(x,t)$ 计算 $h,Y,Z$。作为可容许弱解，$U$ 须满足不等式\eqref{align:ch04-熵积分0}，
而 $\bar{U}$ 作为经典解将以等式满足\eqref{align:ch04-熵积分0}。于是我们推得
\begin{align}
\int_{0}^{T}\int_{\mathbb{R}^{m}}[\partial_{t}\psi h(U,\bar{U})+\sum_{\alpha=1}^{m}\partial_{\alpha}\psi Y_{\alpha}(U,\bar{U})]\,\mathrm{d}x\,\mathrm{d}t+\int_{\mathbb{R}^{m}}\psi(x,0)h(U_{0}(x),\bar{U}_{0}(x))\,\mathrm{d}x
&\geqslant-\int_{0}^{T}\int_{\mathbb{R}^{m}}\Big\{\partial_{t}\psi\theta_{W}(W(\bar{U}))[W(U)-W(\bar{U})]\nonumber\\
&\quad+\sum_{\alpha=1}^{m}\partial_{\alpha}\psi\{\theta_{W}(W(\bar{U}))[X_{\alpha}(U)-X_{\alpha}(\bar{U})]\nonumber\\
&\quad-[\theta_{W}(W(\bar{U}))\Omega(\bar{U})^{\top}-\Xi(\bar{U})^{\top}]M_{\alpha}[U-\bar{U}]\Big\}\,\mathrm{d}x\,\mathrm{d}t\nonumber\\
&\quad-\int_{\mathbb{R}^{m}}\psi(x,0)\theta_{W}(W(\bar{U}_{0}(x)))[W(U_{0}(x))-W(\bar{U}_{0}(x))]\,\mathrm{d}x.\label{align:ch05-poly-ineq1}
\end{align}

接下来对 $U$ 与 $\bar{U}$ 分别写出\eqref{align:ch05-mild-def}，以试验函数 $V^{T}=\psi\theta_{W}(W(\bar{U}))$。这给出
\begin{align}
\int_{0}^{T}\int_{\mathbb{R}^{m}}\Big\{\partial_{t}[\psi\theta_{W}(W(\bar{U}))][W(U)-W(\bar{U})]+\sum_{\alpha=1}^{m}\partial_{\alpha}[\psi\theta_{W}(W(\bar{U}))][X_{\alpha}(U)-X_{\alpha}(\bar{U})]\Big\}\,\mathrm{d}x\,\mathrm{d}t+\int_{\mathbb{R}^{m}}\psi(x,0)\theta_{W}(W(\bar{U}_{0}(x)))[W(U_{0}(x))-W(\bar{U}_{0}(x))]\,\mathrm{d}x=0.\label{align:ch05-poly-id1}
\end{align}
此外，由于 $U$ 与 $\bar{U}$ 都满足对合\eqref{align:ch05-involution}，
\begin{align}
\int_{0}^{T}\int_{\mathbb{R}^{m}}\sum_{\alpha=1}^{m}\partial_{\alpha}\Big\{\psi[\theta_{W}(W(\bar{U}))\Omega(\bar{U})^{\top}-\Xi(\bar{U})^{\top}]\Big\}M_{\alpha}[U-\bar{U}]\,\mathrm{d}x\,\mathrm{d}t=0.\label{align:ch05-poly-id2}
\end{align}

由\eqref{align:ch05-supplementary-matrix}与 $\sum M_{\alpha}\partial_{\alpha}\bar{U}=0$，
\begin{align}
\partial_{t}\theta_{W}(W(\bar{U}))=\partial_{t}W(\bar{U})^{\top}\theta_{WW}(W(\bar{U}))
&=\sum_{\alpha=1}^{m}\partial_{\alpha}X_{\alpha}(\bar{U})^{\top}\theta_{WW}(W(\bar{U}))\nonumber\\
&=-\sum_{\alpha=1}^{m}\partial_{\alpha}\bar{U}^{\top}\mathrm{D}X_{\alpha}(\bar{U})^{\top}\theta_{WW}(W(\bar{U}))\nonumber\\
&=-\sum_{\alpha=1}^{m}\partial_{\alpha}\bar{U}^{\top}[\mathrm{D}W(\bar{U})\mathrm{D}G_{\alpha}(\bar{U})+\Omega(\bar{U})^{\top}M_{\alpha}]^{\top}\theta_{WW}(W(\bar{U}))\nonumber\\
&=-\sum_{\alpha=1}^{m}\partial_{\alpha}\bar{U}^{\top}\mathrm{D}G_{\alpha}(\bar{U})^{\top}\mathrm{D}W(\bar{U})^{\top}\theta_{WW}(W(\bar{U})).\label{align:ch05-poly-Dtheta}
\end{align}
类似地，
\begin{align}
\partial_{\alpha}\theta_{W}(W(\bar{U}))=\partial_{\alpha}\bar{U}^{\top}\mathrm{D}W(\bar{U})^{\top}\theta_{WW}(W(\bar{U})),\label{align:ch05-poly-Dtheta-alpha}
\end{align}
\begin{align}
\partial_{\alpha}[\theta_{W}(W(\bar{U}))\Omega(\bar{U})^{\top}-\Xi(\bar{U})^{\top}]=\partial_{\alpha}\bar{U}^{\top}\mathrm{D}W(\bar{U})^{\top}\theta_{WW}(W(\bar{U}))\Omega(\bar{U})^{\top}-\partial_{\alpha}\bar{U}^{\top}\Gamma(\bar{U})^{\top}.\label{align:ch05-poly-DGamma}
\end{align}
因此，回忆\eqref{align:ch05-Zalpha-polyconvex}，
\begin{align}
\partial_{t}\theta_{W}(W(\bar{U}))[W(U)-W(\bar{U})]+\sum_{\alpha=1}^{m}\partial_{\alpha}\theta_{W}(W(\bar{U}))[X_{\alpha}(U)-X_{\alpha}(\bar{U})]-\sum_{\alpha=1}^{m}\partial_{\alpha}[\theta_{W}(W(\bar{U}))\Omega(\bar{U})^{\top}-\Xi(\bar{U})^{\top}]M_{\alpha}[U-\bar{U}]=\sum_{\alpha=1}^{m}\partial_{\alpha}\bar{U}^{\top}Z_{\alpha}(U,\bar{U}).\label{align:ch05-poly-Z-identity}
\end{align}
由\eqref{align:ch05-DZ-at-bar}，
\begin{align}
\sum_{\alpha=1}^{m}\partial_{\alpha}\bar{U}^{\top}\mathrm{D}Z_{\alpha}(\bar{U},\bar{U})=\left[\sum_{\alpha=1}^{m}M_{\alpha}\partial_{\alpha}\bar{U}\right]^{\top}\Gamma(\bar{U})=0.\label{align:ch05-poly-Z-null}
\end{align}
因此，\eqref{align:ch05-poly-Z-identity}右端关于 $U-\bar{U}$ 为二阶量。

合并\eqref{align:ch05-poly-ineq1}、\eqref{align:ch05-poly-id1}、\eqref{align:ch05-poly-id2}与\eqref{align:ch05-poly-Z-identity}，我们恢复\eqref{align:ch05-relative-ineq2}。证明的剩余部分沿着\cref{theorem:ch05-相对熵稳定性}证明的路线：从\eqref{align:ch05-relative-ineq2}出发，固定任意 $t\in(0,T)$，对 $\sigma\in(0,t)$ 导出\eqref{align:ch05-cutoff-integral}，然后对任意 $\eta(U(\cdot,\tau))$ 的 $L^{\infty}$ 弱$^{*}$ 连续点 $\sigma$ 导出\eqref{align:ch05-cutoff-x}。
这又对由\eqref{align:ch05-u-def}定义的 $u$ 给出\eqref{align:ch05-after-cutoff}，从而\eqref{align:ch05-polyconvex-stability-est}。证明完成。
\end{proof}

特别地，\cref{theorem:ch05-polyconvex-existence}与\cref{theorem:ch05-polyconvex-stability}适用于装备有对合并配备凸偶然熵 $\eta(U)$（取 $W(U)\equiv U$）的守恒律方程组类。人们可尝试通过如下程序将装备有对合并配备多凸偶然熵的更一般方程组类约化到上述特殊类。假设方程组\eqref{align:ch05-cauchy-system}装备有对合\eqref{align:ch05-involution}，并装备有相对于偶然熵 $W$ 多凸\eqref{align:ch05-polyconvex-def}的主偶然熵--熵通量对 $(\eta(U),Q(U))$。我们寻求定义在 $\mathbb{R}^{N}$ 上、分别取值于 $\mathbb{M}^{N\times m}$ 与 $\mathbb{M}^{1\times m}$ 的函数 $S(\Psi)$ 与 $\Pi(\Psi)$，使得
\begin{align}
S(W(U))=W(U),\qquad\Pi(W(U))=Q(U)\label{align:ch05-reduction}
\end{align}
且此外 $(\theta(\Psi),\Pi(\Psi))$ 是\emph{延拓方程组}
\begin{align}
\partial_{t}\Psi(x,t)+\operatorname{div}S(\Psi(x,t))=0\label{align:ch05-extended-system}
\end{align}
的（一般为偶然）熵--熵通量对。

当能找到满足上述规格的函数时，人们可通过先解带初始条件
\begin{align}
\Psi(x,0)=W(U_{0}(x))\label{align:ch05-extended-data}
\end{align}
的\eqref{align:ch05-extended-system}，再由方程 $W(U)=\Psi$ 得到 $U$，来构造 Cauchy 问题\eqref{align:ch05-cauchy-system}、\eqref{align:ch05-cauchy-data}的解。此途径的优点在于\eqref{align:ch05-extended-system}现配备有凸（可能偶然）熵 $\theta$。

上述程序已对弹性动力学与电动力学方程组成功实施。

在弹性动力学中，$U=(F,v)^{\top}$，$\Psi=(F,v,\Theta,\omega)^{\top}$，$\sigma=\sigma(F,\Theta,\omega)$，延拓方程组写为
\begin{align}
\begin{cases}
\partial_{t}F_{i\alpha}-\partial_{\alpha}v_{i}=0,\\
\partial_{t}v_{i}-\partial_{\alpha}\left(\dfrac{\partial\sigma}{\partial F_{i\alpha}}+\dfrac{\partial\sigma}{\partial\Theta_{\beta j}}\dfrac{\partial F_{\beta j}^{*}}{\partial F_{i\alpha}}+\dfrac{\partial\sigma}{\partial\omega}\dfrac{\partial\det F}{\partial F_{i\alpha}}\right)=0,\\
\partial_{t}\Theta_{\beta i}-\partial_{\alpha}\left(\dfrac{\partial F_{\beta i}^{*}}{\partial F_{j\alpha}}v_{j}\right)=0,\\
\partial_{t}\omega-\partial_{\alpha}\left(\dfrac{\partial\det F}{\partial F_{j\alpha}}v_{j}\right)=0,
\end{cases}\label{align:ch05-elastic-extended}
\qquad\begin{aligned}
\alpha&=1,2,3;\ i=1,2,3\\
i&=1,2,3\\
\beta&=1,2,3;\ i=1,2,3
\end{aligned}
\end{align}
而熵--熵通量对为
\begin{align}
\theta=\frac{1}{2}|v|^{2}+\sigma(F,\Theta,\omega),\label{align:ch05-elastic-theta}
\end{align}
\begin{align}
\Pi_{\alpha}=-\left(\frac{\partial\sigma}{\partial F_{i\alpha}}+\frac{\partial\sigma}{\partial\Theta_{\beta j}}\frac{\partial F_{\beta j}^{*}}{\partial F_{i\alpha}}+\frac{\partial\sigma}{\partial\omega}\frac{\partial\det F}{\partial F_{i\alpha}}\right)v_{i}.\label{align:ch05-elastic-Pi}
\end{align}
在"流形" $\Psi=W(U)=(F,v,F^{*},\det F)^{\top}$ 上，\eqref{align:ch05-elastic-extended}约化为方程组\eqref{align:ch03-等熵热弹性系统}（$b=0$）连同运动学守恒律\eqref{align:ch03-余因子守恒}、\eqref{align:ch03-行列式守恒}，


而 $(\theta,\Pi)$ 约化为\cref{section:ch03-例子}记录的经典熵--熵通量对。

在电动力学中，对 Born--Infeld 本构关系，$U=(B,D)^{\top}$，$\Psi=(B,D,P)^{\top}$，延拓方程组写为
\begin{align}
\begin{cases}
\partial_{t}B+\operatorname{curl}[\theta^{-1}(D+B\wedge P)]=0,\\
\partial_{t}D-\operatorname{curl}[\theta^{-1}(B-D\wedge P)]=0,\\
\partial_{t}P-\operatorname{div}[\theta^{-1}(I+BB^{\top}+DD^{\top}-PP^{\top})]=0,
\end{cases}\label{align:ch05-em-extended}
\end{align}
而熵--熵通量对为
\begin{align}
\theta=(1+|B|^{2}+|D|^{2}+|P|^{2})^{\frac{1}{2}},\label{align:ch05-em-theta}
\end{align}
\begin{align}
\Pi=P-\theta^{-2}[P-D\lambda B-(D\cdot P)D-(B\cdot P)B].\label{align:ch05-em-Pi}
\end{align}
同样，在"流形" $\Psi=W(U)=(B,D,D\wedge B)^{\top}$ 上，\eqref{align:ch05-em-extended}约化为 Maxwell 方程组\eqref{align:ch03-Maxwell}（$J=0$）连同补充守恒律\eqref{align:ch03-动量电磁}，

而 $(\theta,\Pi)$ 约化为\eqref{align:ch03-BornInfeld}中记录的熵--熵通量对 $(\eta,Q)$。

\end{document}
